\section{Stack tecnológico}\label{sec:stack}
Las herramientas se eligieron por su pertinencia para el problema, su instalación libre de licencias (reproducibilidad) y su integración con Python.

\begin{table}[H]
\centering\small
\caption{Stack tecnológico.}
\begin{tabularx}{\textwidth}{L{3.2cm}L{4.6cm}Y}
\toprule
Tecnología & Función & Justificación \\
\midrule
Python 3 & Lenguaje de implementación (paquete \texttt{bus\_sched}) & Lenguaje de referencia del curso; LS y LPT resuelven cada ruta en milisegundos. \\
Git / GitHub & Control de versiones & Historial del trabajo y de las contribuciones de cada integrante. \\
CSV / JSON & Instancias, perfiles de tráfico y resultados & Formatos abiertos, legibles y reproducibles. \\
openpyxl / pandas & Lectura del Excel del proyecto y análisis de resultados & Permiten leer los parámetros, validar los datos y generar tablas. \\
matplotlib & Gráficos y diagramas de Gantt & Figuras reproducibles de los resultados. \\
pytest & Pruebas automatizadas & Casos parametrizados y verificación continua. \\
Jupyter Notebook & Demostración del prototipo & Ejecución paso a paso (también en Google Colab). \\
\LaTeX & Informe & Documento reproducible; las tablas de resultados se generan desde el programa. \\
OR-Tools (CP-SAT) & Prevista para la referencia exacta posterior & Libre; modela intervalos sin solapamiento y duraciones que dependen del inicio. \\
\bottomrule
\end{tabularx}
\end{table}
